\documentclass[11pt,a4paper]{article}
\usepackage[T1]{fontenc}
\usepackage[utf8]{inputenc}
\usepackage{lmodern}
\usepackage{amsmath,amssymb,amsthm,mathtools,mathrsfs}
\usepackage{graphicx,float,xcolor,multicol}
\usepackage[shortlabels]{enumitem}
\usepackage[margin=27mm]{geometry}
\usepackage{microtype,needspace,placeins}
\usepackage{tikz,tikz-cd}
\usetikzlibrary{arrows.meta,calc,positioning,patterns,decorations.pathreplacing}
\usepackage[colorlinks=true,linkcolor=teal,urlcolor=teal,citecolor=teal]{hyperref}
\setlength{\parindent}{0pt}
\setlength{\parskip}{.6em}
\setcounter{secnumdepth}{0}
\allowdisplaybreaks
\raggedbottom
\providecommand{\cross}{\times}
\DeclareUnicodeCharacter{2020}{\ensuremath{\dagger}}
\DeclareMathOperator{\supp}{supp}
\DeclareMathOperator*{\esssup}{ess\,sup}
\providecommand{\chapter}{\section}
\newenvironment{tool}[2]{\subsection*{Tool #1 --- #2}}{}
\newcommand{\ToolRef}[1]{Tool #1}
\newcommand{\FigureTag}[2]{\textit{Figure #1}}
\newcommand{\Result}[1]{\subsection*{#1}}
\newcommand{\Application}[1]{\subsubsection*{Application\if\relax\detokenize{#1}\relax\else\space--- #1\fi}}
\newcommand{\ProofHeading}{\par\textit{Proof.}\space}
\newcommand{\ProofOf}[1]{\par\textit{Proof of #1.}\space}
\newcommand{\RemarkHeading}{\par\textbf{Remark.}\space}
\newcommand{\NoteHeading}{\par\textbf{Note.}\space}
\newcommand{\ExerciseHeading}{\par\textbf{Exercise.}\space}

\graphicspath{{figures/surgeries/}}
\begin{document}
\section*{Sullivan's No Wandering Theorem}
% BLOG-CONTENT-BEGIN

\subsection*{Theorem 4.1 (No Wandering Theorem)}
Rational maps of degree at least $2$ on $\hat{\mathbb{C}}$ have no wandering Fatou components.

        A Fatou component of $f$ may be invariant, periodic, preperiodic, or wandering. Sullivan's theorem rules out the wandering case for rational maps; see \cite{branner2014quasiconformal,sullivan1985quasiconformal}. We begin with the analytic observations needed for the surgery proof. 

        Let $U$ be a wandering Fatou component, and let $f^{\circ n}(U) = U_n$, then $U_n \cap U_m = \varnothing$ $(\forall n \neq m)$, and since $f$ has finitely many critical points, eventually, $U_n$ will have no critical points of $f$. Conjugate $f$ by a M\"obius transform (if need be) to assume WLOG that $\infty \in f^{-1}(U)$. That is, $\infty$ belongs to a Fatou component in $f^{-1}(U)$. As all $U_n$ are disjoint, they won't intersect the component that contains $\infty$ in $f^{-1}(U)$. This forces $\bigcup_{n\geq 1}U_n$ to be bounded in $\mathbb{C}$ (after stereographic projection), as the domains $U_n$ won't cluster around $\infty$ in $\hat{\mathbb{C}}$.

        Let $g$ be a non-constant limit function of the family $\{f^{\circ n}\big|_U\}_n$. Then there exists $z_0 \in U$ such that $g'(z_0) \neq 0$. Let $\overline{B_{\delta}(z_0)} \subset U$ (closed ball of radius $\delta$) be such that $g$ is injective on it. Set,
        $$\epsilon = \inf_{z\in \partial B_{\delta}(z_0)}{|g(z)-g(z_0)|}$$
        
        Let $f^{\circ n_j}\rightarrow g$ uniformly on compact subsets of $U$. Going by first principles, $\exists J$ such that $\forall j$, $j\geq J$, $|f^{\circ n_j}(z)-g(z)| < \epsilon$ $(\forall z \in \overline{B_{\delta}(z_0)})$. Set $F_j(z):= f^{\circ n_j}(z)-g(z_0)$ and $G(z):= g(z)-g(z_0)$. Then $\forall j > J$, and $z\in \partial B_{\delta}(z_0)$,
        $$|F_j(z) -G(z)| = |f^{\circ n_j}(z)-g(z)|<\epsilon \leq |G(z)|$$
        
        $\implies$ $G(z)$ and $F_j(z)$ have the same number of zeroes in $B_{\delta}(z_0)$ (by Rouch\'e's Theorem). So, $\forall j \geq J$, $\exists$ $z^*(j)$ such that $f^{\circ n_j}(z^*(j)) = g(z_0)$. This gives,
        $$g(z_0) \in \bigcap_{j\geq J}{U_{n_j}}$$
        
        $\implies$ $U$ isn't a wandering domain $\implies$ $g$ is constant.

        This forces $\operatorname{diam}(f^{\circ n}(L)) \rightarrow 0$ for every compact $L\Subset U$. If not, then $\exists L$, $L \subset U$, compact in $U$, and a subsequence $\{n_j\}_j$ such that diam$(f^{\circ n_j}(L)) \geq \delta$. Normality yields another subsequence which we shall again denote by $\{n_j\}_j$ such that $f^{\circ n_j} \rightarrow (z\rightarrow c)$ (the constant function $c$). This implies $f^{\circ n_j}(L) \subset B_{\epsilon}(c)$ $(\forall \epsilon >0)$. Contradiction.

% BLOG-CONTENT-END
\bibliographystyle{amsalpha}
\bibliography{tools/profiles/surgeries/MSc_Dissertation}
\end{document}
